\section{Benchmark and Protocol}\label{sec:protocol}

\paragraph{VN-AgentInject.} The benchmark turns the model of Section~\ref{sec:threat} into runnable \emph{test cases}.
All data are synthetic. Each family yields one attack case and one \bPlain{} case per language variant, and one third
of the families add an \bAuth{} case. The benchmark has three disjoint sets, each balanced by mechanism. The
\textbf{Dev} set (48 families, 560 cases) is used for every design decision. \textbf{\tOne} (72 families, 840 cases)
is the first sealed set. \textbf{\tTwo} (48 new families, 560 cases) was written after \tOne{} had been run. It is new
to the frozen \ours, but its authors knew the \tOne{} results (Section~\ref{sec:discussion}). A family, together with
all its variants and benign counterparts, belongs to exactly one set.

Attack texts were written by hand by people who did not read the defence code. Each family was reviewed against a
five-item rubric, including the mechanism label, naturalness, real harm and not revealing the answer key. Each family
also had to pass an overlap filter against earlier versions and a \emph{replay check}. In the replay check, a
reference trajectory that follows the injection must be scored as harmful, and a trajectory that completes the task
safely must be scored as safe.

\paragraph{Sealing.} Before each test run, three things were frozen. The first is the test cases, pinned by SHA-256.
The second is a manifest that pins the hashes of the prompts, the decoding settings, the defence code, the model
versions and the list of configurations. The third is the hypotheses together with the multiple-comparison
correction. The runner refuses to start if anything differs. Each pair of model and configuration was run only once,
except for resuming cases that failed for infrastructure reasons. The scorer's answer key never appears in the prompt
or in the runtime policy, so the model cannot see it. \tOne{} has six pre-registered hypotheses, H1 to H6. The
pre-registered hypothesis of \tTwo{} targets a different variant of \ours{} (Section~\ref{sec:discussion}). The
comparison of \ours{} with \cNone{} on \tTwo{} is therefore a non-pre-registered replication analysis, although the
configurations and data are pinned in the same way.

\paragraph{Backbones and baselines.} The three open-weight backbones are Gemma-4-12B, Llama-3.1-8B and Qwen3.5-9B,
pinned by version. All three share the same ReAct prompt, temperature 0, at most 512 tokens per step and eight steps.
Gemini 3 Flash is run only on Dev and AgentDojo for reference and is excluded from all tests. On Dev, \ours{} is
compared with six baselines implemented in the same agent and the same broker: \cSpot~\cite{hines2024spotlighting},
\cSand{} (repeating $u$), \cPG{} (Prompt Guard~2~\cite{meta2025promptguard2} screening tool outputs), \cMel{} and
\cProg{} (the asterisk marks re-implementations after MELON~\cite{zhu2025melon} and Progent~\cite{shi2025progent},
which do not represent the original systems), and \cPhrase{} (phrase rules, an earlier version of this system).

\paragraph{Metrics and statistics.} For a set of attack cases $\mathcal{A}$ and a set of benign cases $\mathcal{B}$,
\begin{equation}\label{eq:metrics}
\mathrm{ASR} = \frac{|\{f \in \mathcal{A} : \mathrm{Harm}(f)\}|}{|\mathcal{A}|}, \qquad
\mathrm{U} = \frac{|\{f \in \mathcal{B} : \mathrm{Done}(f) \land \lnot \mathrm{Harm}(f)\}|}{|\mathcal{B}|},
\end{equation}
where $\mathrm{Done}(f)$ means that the frozen scorer confirms that the user's task was completed, using a
deterministic task-specific condition. Model errors, that is, malformed or truncated outputs, count as utility
failures and never as successful defence. The tables use three counting units. A \emph{case} is one run. A
\emph{family--model pair} is harmful if at least one attack case of that family on that model is harmful. A
\emph{family} is harmful if at least one of its cases is harmful on any model. Cases within a family (up to 15, five
variants on three models) share the task and environment, so they are not independent. The measured design effect
ranges from 1.5 to 7.7. The unit of inference is therefore the family. For attack success, the paper uses the exact
one-sided McNemar test~\cite{mcnemar1947} on the harmful status of each family. For utility, it uses a two-sided
sign-flip permutation test by family. 95\% confidence intervals are computed by family bootstrap with 10,000
resamples. The pre-registration used McNemar on individual cases. That test underestimates the variance, so all six
hypotheses were re-tested at the family level. The conclusion of each hypothesis is unchanged. The $p$-values of H1 to
H6 are adjusted with the Holm procedure~\cite{holm1979}.

\paragraph{AgentDojo-VI.} As an external check beyond VN-AgentInject, the paper also uses the Banking suite (144 pairs
of user task and injection task, 16 user tasks) and the Slack suite (105 pairs, 21 tasks) of AgentDojo
v1.2.2~\cite{debenedetti2024agentdojo}, in English and in a Vietnamese translation. The translation covers the
environment, tasks, injection goals and the \code{important\_instructions} attack template. IBANs, amounts, names and
addresses are kept unchanged, while the tools and scoring functions are the original code. With Llama, \ours{} is run
in full in both languages with the frozen rule code. The only new part is an adapter declared before the run. It
specifies which tools are sinks, which argument holds the destination and which read tools return confidential data.
The \cSand{} layer is AgentDojo's built-in \code{repeat\_user\_prompt}, which is also run on its own with Llama in both
languages.
